\section{Experiments: Predictions and Scope}\label{sec:scope}
\begin{table}[t]
\centering\small
\setlength{\tabcolsep}{3.4pt}
\begin{tabular}{@{}lrrr@{}}
\toprule
25-call prefix & Class 1 & All classes & Squared error\\
\midrule
Restricted switching & 0 & 0 & $.322$\\
Shared geometric & 23 & 0 & $.249$\\
Shared platform & 23 & 0 & $.249$\\
Observation preset & 23 & 23 & $1.55\cdot10^{-12}$\\
\bottomrule
\end{tabular}
\caption{Frozen scale-selection diagnostic on 23 certified models.
The platform rows isolate allocation at the preset's selected center.
Entry columns count models where all sampled draws for the named
classes lie in their terminal rectangles at the switch. Squared error
is prior-weighted within model and averaged across models.
Each policy uses the same 106 calls.}
\label{tab:selected_scales}
\end{table}

\paragraph{Choosing scales before running the solver.}
We fix a $5\times5$ grid of heights and masses, with
$a=1$, $r=1.29$, $\alpha=.001$, and $\tau=\nu=\alpha^2$.
For each model we compute the largest certified low cutoff
on the $.01$ grid from $.08$ to $.30$ and the smallest
certified high cutoff on the $.1$ grid from $1.5$ to $8$.
We screen width-$.02$ bands with outward-rounded evaluations
of the barrier, full channel, other-class initialization, and
terminal tests. The first-pass rule selected the same center,
$.40$, on all 23 passing models. The second rule evaluates
centers from $.24$ to $1.10$ at spacing $.02$ and minimizes
the certified first-entry bound, breaking ties by smaller
center (Supplement S28). The rule uses the prior geometry
and prescribed noise tolerances. We freeze it before drawing
32 new independent pairs of component and measurement noise
per class and model.

Figure~\ref{fig:selected_scales} shows the entire selection
grid. The high cutoff varies from $1.5$ to $4.0$. At $h=1.78$,
heights $1.6,2,2.4$ yield centers $.36,.44,.60$ and
certified bounds $24,17,14$, respectively.
For $d=1.6$ with $h=2.1$ or $2.4$, the channel tests certify
no candidate band within 25 calls on the search grid.
Across the 23 certified models, all 2208 observations fall in $E$.
The observed class-1 first entries under the preset take
14--18 calls. The certificate ranges from 14 to 24 calls.
The difference between the bound and the largest observed
entry call is 0--6, with a median of one call and equality
in four models. The 25-call prefix matches the comparison
budget and the prepared-strip test covers the remaining calls.

We compare restricted switching, a geometric prefix, a shared
platform at the selected center, and the $0/25/0$ preset.
All use a 25-call switch window and 106 total calls on paired
observations. At the preset's selected center, allocating its
intermediate calls to every observation leaves class 0 outside $G_0$
at the switch (Table~\ref{tab:selected_scales}).

\paragraph{Independently tuned shared platforms.}
Shared and gated platforms independently search the same
2275 variance/length pairs on noiseless model observations.
At 106 calls, both yield correct nearest templates on all 5760
new noisy draws over 23 models, with prior-weighted squared
error $1.47\cdot10^{-12}$. Both use 11--14 prefix calls.
On the nominal model, a 25-call shared platform at $v=.25$
recovers all 1536 draws despite switching outside $G_0,G_1$;
the preset's center $.44$ yields 1024 correct draws when shared.
On the noiseless grid, gating admits 428--1142 passing pairs
per model versus 9--25 for sharing (Supplement S28).

On the prespecified nominal fixture, holding the high endpoint
for 12 calls and the low endpoint for 13 skips the candidate band.
Replacing the first low call by one band call visits the band
but still leaves class 1 at $s<.194$ at the switch. Neither prefix
produces any class-1 terminal entries in 106 calls. Holding the
selected center for the full prefix reaches $G_1$ at call 16 for
all 32 class-1 draws. The shared geometric schedule reaches it at
call 19 even though its grid misses this narrow certified band;
the band gives a sufficient guarantee for the specified prefix.

\paragraph{Dependence on the solver.}
In the matched fine-calibration diagnostic, terminal-only RED
from a least-squares start gives $1.47\cdot10^{-20}$ squared error
in 154 calls, comparable to the preset from backprojection.
The direct observation-based template estimator is also consistent.
Changing the initial state or the data and regularization forces
can remove the common barrier. Supplement S21 records the
initialization, coefficient, and precision diagnostics.

\paragraph{Executing the certified budget.}
The analytic cap of 14--24 calls yields entry on all 23 models.
Prior-weighted intermediate calls fall from $8.05$ to $5.54$
per observation ($31\%$). Errors match to reported precision
at 106 calls and with only 81 post-prefix calls (95--105 total,
mean $98.2$; Supplement S28).

\paragraph{When schedule history disappears.}
We fit eight common-isotropic components to 2048 grayscale
$4\times4$ patches and use periodic blur and Gaussian noise
of standard deviation $.01$. Condition~\eqref{eq:ind_coeff}
implies $\sigma_L/\tau\ge\sqrt{333}\approx18.25$ in the
risk theorem. In the frozen patch model the ratio is $.208$,
and changing $\sigma$ from $.1$ to $.02$ changes $c(v)$ from
$.0520$ to $.00415$, a factor of about $12.5$.
Holding $c$ at its terminal value while varying the posterior
weights reduces the original flower diagnostic's 154-call
mean squared output difference from $1.31\cdot10^{-5}$ to
$1.49\cdot10^{-9}$. This modified update isolates the effect
of changing weights (Supplement S23).

In the frozen four-photo validation, the intermediate protocol
changes 154-call squared error by $-.138\%,+.079\%,+.730\%$
and $-.377\%$ relative to terminal only. Whole-ball covariance
bounds certify that the three exported 154-call states share
a contracting terminal region on each of the 3072 observations.
These are local certificates constructed using later states.
They cover subsequent exact terminal maps and the measured
schedule differences decay (Supplement S25). On a selected
development patch, two fixed-scale prefixes enter a higher-error
region while low/high switching reaches the intermediate
protocol's lower-error region at the same total budget
(Supplement S26). This shows why a bound for arbitrary switching
needs a common invariant region.

\section{Conclusion}
For a fixed RED solver initialized by backprojection, a common
invariant region can prevent recovery under every restricted scale
policy, regardless of its finite budget. Intermediate scales can
open a channel that changes the limiting recovery risk. The
geometric tests certify this mechanism and predict entrance budgets
on the structured family. Independently tuned shared platforms
show that successful allocation need not be unique; the photo tests
identify a regime where terminal contraction erases schedule history.
Extending the positive risk separation to broader sensing geometries
and image priors remains open.
